\section{Related work}
\label{sec:lit}

\paragraph{Procurement transport with inventory decisions.}
\autocite{Berman2006InboundCost} and
\autocite{Sancak2011Multi-itemPolicy} price full and
less-than-truckload service against inventory under a volumetric
capacity; \autocite{Schoneberg2010AnNetworks} and
\autocite{Schoneberg2013ANetworks} do so for multi-supplier inbound
networks; \autocite{Diaz-Madronero2014APlanning} formulate the
operational single-supplier full-truckload problem that the benchmark
of Appendix~\ref{appendix:A} extends, with capacity as a maximum
container count; \autocite{Hein2016QuantitativeProblem},
\autocite{Diaz-Madronero2017ADecisions},
\autocite{Meyer2018TransportIndustry}, \autocite{Chitsaz2019ARouting},
\autocite{Baller2022OptimizingApproach} and
\autocite{Hashemi-Amiri2023IntegratedApproach} extend the setting to
milk-runs, routing and uncertainty. In every case the trailer is a
scalar: a volume, a weight, or a number of units. Table
\ref{table:ComProTra} positions these models on the features that
matter here, and in no prior row is the capacity a condition on
placement.

\paragraph{RTIs.}
The closed loop itself is modelled at a tactical or network level.
\autocite{SOYSAL2023101504} plan vendor-managed inventory of agro-food
crates with forward and reverse flows; \autocite{GUO2024104339}
optimise a two-tier container management network with return penalties;
\autocite{Zhang17012024} formulate a perishable inventory-routing
problem with several RTI types. These carry the reverse flow but
represent the vehicle as a capacity, so a folded container's occupying
a different volume from an erected one changes a coefficient and not
the structure of the problem.

\paragraph{Container loading.}
The loading side is surveyed by the typology of
\autocite{Wascher2007AnProblems}, the review of practically relevant
constraints of \autocite{Bortfeldt2013ConstraintsReview} and the
exact-method survey of \autocite{Silva2019ExactProblems}. Our loading
requirements --- orientation, stacking with load bearing, weight limits
--- are those studied as side constraints of a single-container problem
by \autocite{Junqueira2012Three-dimensionalConstraints} and
\autocite{Alonso2020AConstraints}. The nearest methodological
precedent is \autocite{Mahvash2018}, who solve three-dimensional bin
packing by column generation in which a column is a packing
configuration of one bin: we take the same unit of decision, but the
master it feeds is a multi-period transport and inventory problem
rather than a set-covering model, so the columns of a direction are
generated once and priced against duals that move with the inventory
balance.

% Table 1 of main.tex (table:ComProTra), re-set one row per line.  The data are
% those of main.tex:139-284; the only cell changed is 3D-SCPTOP's own RTI entry,
% which was blank there and is M|F (several types, foldable).
\begin{table}[!htbp]
  \centering
  \small
  \caption{Procurement transport models with inventory decisions, on the
  features that separate them from 3D-SCPTOP. \emph{Config.}: SS single
  source and single destination, MS multi-supplier, MM multi-echelon.
  \emph{Level}: O operational, T tactical. \emph{Mode}: FTL
  full-truckload, LTL less-than-truckload, MR milk-run, CES courier
  express. \emph{Capacity}: V volumetric, W weight, UH a count of
  loaded units. \emph{RTIs}: S one type, M several, F foldable, NF not.
  \emph{Framework}: (F)(MO)(I)LP linear programming variants, CCP
  chance-constrained. A dash is a feature the model does not carry.}
  \label{table:ComProTra}
  \begin{adjustbox}{width=\textwidth}
  \begin{tabular}{llllllllll}
    \toprule
    Article & Config. & Level & Mode & Capacity & Tr.\ cost & Inv.\ cap. & RTIs & Inv.\ cost / safety stock & Framework \\
    \midrule
    \autocite{Berman2006InboundCost}            & MM & T   & FTL\,$|$\,LTL            & V        & X & ---  & ---      & X / ---  & MINLP  \\
    \autocite{Sancak2011Multi-itemPolicy}       & ---& O   & LTL\,$|$\,FTL            & V        & X & ---  & ---      & X / X    & MILP   \\
    \autocite{Schoneberg2010AnNetworks}         & MS & T   & FTL\,$|$\,LTL            & V        & X & X    & ---      & X / ---  & MILP   \\
    \autocite{Schoneberg2013ANetworks}          & MS & O   & FTL\,$|$\,LTL            & V        & X & X    & ---      & X / ---  & MILP   \\
    \autocite{Diaz-Madronero2014APlanning}      & SS & O   & FTL                      & UH       & X & ---  & ---      & X / X    & FMOILP \\
    \autocite{Hein2016QuantitativeProblem}      & MS & O   & MR                       & V        & X & ---  & ---      & X / ---  & MILP   \\
    \autocite{Diaz-Madronero2017ADecisions}     & MS & T   & FTL\,$|$\,LTL\,$|$\,MR   & V        & X & X    & ---      & X / X    & MILP   \\
    \autocite{Meyer2018TransportIndustry}       & MS & T, O& MR\,$|$\,FTL             & V        & X & X    & ---      & X / ---  & MILP   \\
    \autocite{Chitsaz2019ARouting}              & MS & O, T& MR                       & V        & X & X    & ---      & X / ---  & MILP   \\
    \autocite{Baller2022OptimizingApproach}     & MS & O   & MR\,$|$\,FTL\,$|$\,CES   & V\,$|$\,W& X & X    & ---      & --- / X  & MILP   \\
    \autocite{Hashemi-Amiri2023IntegratedApproach}& MS & T & MR                       & V        & X & ---  & ---      & --- / --- & CCP   \\
    \autocite{GUO2024104339}                    & MM & T, O& FTL                      & UH       & X & X    & S\,$|$\,NF & X / --- & MINLP \\
    \autocite{Zhang17012024}                    & MS & T, O& MR                       & UH       & X & X    & M\,$|$\,NF & X / --- & ILP   \\
    \midrule
    \textbf{3D-SCPTOP}                          & SS & O   & FTL & \textbf{3D loading} & X & X & M\,$|$\,F & X / X & MILP \\
    \bottomrule
  \end{tabular}
  \end{adjustbox}
\end{table}

